% ============================================================
%  Capítulo 4 — Análisis
%  TT 2026-B066 | ESCOM-IPN
%  Frausto Robles Ángel Ali · Rodríguez Verdín Sandoval Vanesa
% ============================================================

\chapter{Análisis}

En este capítulo se definen los requerimientos del sistema a partir de los prototipos
de interfaz desarrollados para el flujo de trabajo del laboratorio colaborador y de
reuniones con el Dr.\ César Augusto Sandino Reyes López, investigador del Laboratorio
de Neurociencia Conductual de la ENMyH-IPN. Se documentan
también las reglas de negocio, los riesgos identificados y la factibilidad del proyecto
desde cuatro ángulos: tecnológico, operativo, económico y legal.

Los requerimientos funcionales se organizan por módulo de la interfaz. Los tiempos de
respuesta y criterios de rendimiento se derivan de los atributos de calidad definidos
en la norma ISO/IEC~25010:2023 \cite{iso25010}: \textit{time behaviour},
\textit{functional correctness} y \textit{usability}. El ciclo de levantamiento sigue
los procesos de definición de requerimientos de la norma ISO/IEC/IEEE~12207:2017
\cite{iso12207}.

% -----------------------------------------------------------
\section{Requerimientos}
% -----------------------------------------------------------

\subsection{Requerimientos Funcionales (RF)}

Los requerimientos funcionales están organizados por módulo de la interfaz con tres
niveles de prioridad: \textbf{Alta} (sin esto el prototipo no funciona),
\textbf{Media} (importante pero puede quedar para una iteración posterior) y
\textbf{Baja} (deseable pero no urgente). Cada requerimiento se derivó directamente
de los prototipos de interfaz diseñados \cite{pressman2014,sommerville2016}.

\begin{longtable}{|>{\bfseries}p{1.1cm}|p{2.0cm}|p{6.7cm}|p{1.3cm}|p{1.7cm}|}
\caption{Requerimientos Funcionales (RF).
\label{tab:rf}}\\
\hline
\textbf{ID} & \textbf{Módulo} & \textbf{Descripción} & \textbf{Prior.} & \textbf{Actor} \\
\hline
\endfirsthead
\hline
\textbf{ID} & \textbf{Módulo} & \textbf{Descripción} & \textbf{Prior.} & \textbf{Actor} \\
\hline
\endhead
\hline
\endfoot

% ── Autenticación ──────────────────────────────────────────
RF-01 & Autenticación
      & El sistema debe permitir iniciar sesión con correo electrónico institucional
        (\texttt{@ipn.mx}) y contraseña. Las credenciales incorrectas muestran
        un mensaje de error sin revelar cuál campo falló.
      & Alta & Investigador / Admin \\
\hline
RF-02 & Autenticación
      & El sistema debe cerrar la sesión del usuario de forma segura, invalidando
        el token JWT activo. La opción de cierre de sesión debe ser accesible desde
        el menú de usuario en todas las pantallas.
      & Alta & Investigador / Admin \\
\hline
RF-03 & Autenticación
      & El sistema debe ofrecer recuperación de contraseña enviando un enlace de
        restablecimiento al correo registrado.
      & Media & Investigador \\
\hline
RF-04 & Autenticación
      & El sistema debe gestionar dos roles con permisos distintos: Investigador
        (acceso a sus propios experimentos) y Administrador (acceso a todo el sistema).
      & Alta & Sistema \\
\hline

% ── Gestión de usuarios ────────────────────────────────────
RF-05 & Gestión de usuarios
      & El Administrador debe poder crear, modificar y desactivar cuentas de usuario
        desde el panel de administración. La creación requiere nombre, correo y
        contraseña temporal; el usuario la cambia al primer ingreso.
      & Alta & Admin \\
\hline
RF-06 & Gestión de usuarios
      & El Investigador debe poder actualizar su nombre, correo y contraseña desde
        su perfil de usuario.
      & Media & Investigador \\
\hline
RF-07 & Gestión de usuarios
      & \textbf{[PENDIENTE — entrevista con el Dr.\ Sandino]:} ¿Los
        investigadores pueden registrarse solos con confirmación por correo, o solo
        el Administrador puede crear cuentas? Definir antes de iniciar
        la implementación del módulo de autenticación.
      & Alta & Admin \\
\hline

% ── Carga de video ─────────────────────────────────────────
RF-08 & Carga de video
      & El sistema solo acepta archivos en formato \texttt{.mp4}. Cualquier otro
        formato se rechaza con un mensaje de error específico antes de intentar
        subirlo al servidor.
      & Alta & Sistema \\
\hline
RF-09 & Carga de video
      & El sistema debe permitir subir uno o dos videos por experimento: Día~1
        (sesión basal, 20~min) y Día~2 (sesión post-tratamiento, 5~min). Ambos son
        opcionales individualmente; el análisis usa los que se proporcionen.
      & Alta & Investigador \\
\hline
RF-10 & Carga de video
      & El sistema debe validar que el archivo subido sea un video reproducible antes
        de encolarlo. Si la validación falla, lo rechaza con código de error y
        mensaje claro.
      & Alta & Sistema \\
\hline
RF-11 & Carga de video
      & Al crear un experimento el investigador debe registrar: nombre del
        experimento, fecha, número esperado de animales (típicamente~4), tratamiento o
        condición experimental, y notas adicionales opcionales.
      & Alta & Investigador \\
\hline
RF-12 & Carga de video
      & El análisis debe iniciarse automáticamente al completarse la carga y
        validación del video, sin intervención del investigador.
      & Alta & Sistema \\
\hline

% ── Análisis conductual ────────────────────────────────────
RF-13 & Análisis conductual
      & El pipeline de análisis ejecuta cuatro etapas secuenciales: preprocesamiento
        del video (decodificación y corrección de contraste), detección de cilindros,
        seguimiento de animales (\textit{tracking}) y clasificación de conductas.
        Cada etapa tiene un indicador de estado en la pantalla de progreso.
      & Alta & Sistema \\
\hline
RF-14 & Análisis conductual
      & El sistema debe detectar automáticamente la posición de los cuatro cilindros
        en el encuadre con una confianza mínima de 0.70. Si no se alcanza ese umbral,
        el pipeline se detiene y genera un reporte de diagnóstico.
      & Alta & Sistema \\
\hline
RF-15 & Análisis conductual
      & Mientras el análisis corre, el investigador debe ver el progreso en tiempo
        real con una barra de porcentaje y la etapa activa. El análisis continúa
        en segundo plano aunque el investigador cierre la pestaña.
      & Alta & Investigador \\
\hline
RF-16 & Análisis conductual
      & Si el análisis falla, el sistema muestra el código de error, la descripción
        del problema y un botón para descargar el reporte de diagnóstico en PDF.
        El investigador no puede reiniciar el análisis desde esta pantalla (ver RN-04).
      & Alta & Investigador \\
\hline
RF-17 & Análisis conductual
      & El sistema debe clasificar cuadro a cuadro la conducta de cada animal: nado
        activo, inmovilidad o escalamiento \cite{armario2021}. Las conductas son
        mutuamente excluyentes (ver RN-07).
      & Alta & Sistema \\
\hline

% ── Resultados ─────────────────────────────────────────────
RF-18 & Resultados
      & La pantalla de resultados debe mostrar el tiempo total en segundos y el
        porcentaje de cada conducta por animal y por sesión, en una tabla
        descargable.
      & Alta & Investigador \\
\hline
RF-19 & Resultados
      & El sistema debe ofrecer un desglose de conductas por minuto para cada animal.
      & Media & Investigador \\
\hline
RF-20 & Resultados
      & Cuando ambos videos del mismo experimento estén procesados, el sistema debe
        mostrar una comparación Día~1 vs.\ Día~2 con barras de distribución por
        animal.
      & Alta & Investigador \\
\hline
RF-21 & Resultados
      & El investigador debe poder descargar el reporte de resultados en PDF y en
        CSV desde la pantalla de resultados y desde el dashboard.
      & Alta & Investigador \\
\hline

% ── Dashboard / gestión de experimentos ──────────────────
RF-22 & Dashboard
      & El dashboard debe mostrar la lista de experimentos del investigador activo
        con nombre, fecha, tratamiento, sesión, estado (en proceso, completado, error)
        y accesos directos a resultados y descarga.
      & Alta & Investigador \\
\hline
RF-23 & Dashboard
      & El sistema debe mostrar un aviso visible cuando un video está a menos de
        7~días de ser eliminado automáticamente, indicando el nombre del experimento
        y la fecha de vencimiento exacta. El período de 7~días da tiempo suficiente
        para descargar reportes antes del borrado, conforme al atributo
        \textit{user error protection} de la ISO/IEC~25010:2023 \cite{iso25010},
        que establece que el sistema debe proteger al usuario de consecuencias
        no intencionadas de sus acciones u omisiones.
      & Media & Sistema \\
\hline
RF-24 & Dashboard
      & El sistema debe borrar automáticamente los archivos de video 30~días después
        de que el análisis termine exitosamente, notificando al investigador con un
        aviso en la interfaz antes de que ocurra (ver RN-05).
      & Media & Sistema \\
\hline
RF-25 & Dashboard
      & Los resultados y reportes deben conservarse indefinidamente, aunque el video
        original ya haya sido borrado.
      & Alta & Sistema \\
\hline
RF-26 & Dashboard
      & El investigador debe poder eliminar un experimento desde el dashboard. La
        acción es permanente e irreversible y requiere confirmación explícita (ver
        RN-09).
      & Media & Investigador \\
\hline

% ── Administración ─────────────────────────────────────────
RF-27 & Administración
      & El Administrador debe poder ver la lista de todos los usuarios con correo,
        rol, estado (activo/inactivo), último acceso y número de experimentos. La
        lista es filtrable por rol y estado.
      & Media & Admin \\
\hline
RF-28 & Administración
      & El Administrador debe poder monitorear el estado del sistema: porcentaje de
        uso del disco, espacio disponible en GB, tamaño de la cola de procesamiento,
        número de análisis activos y ETA estimado. El sistema alerta cuando el disco
        supera el 80\,\% de uso.
      & Baja & Admin \\
\hline
RF-29 & Administración
      & El Administrador puede ver y gestionar los experimentos de cualquier
        investigador (ver RN-08).
      & Media & Admin \\
\hline

\end{longtable}

% -----------------------------------------------------------
\subsection{Requerimientos No Funcionales (RNF)}
% -----------------------------------------------------------

Los tiempos de respuesta se derivan del atributo \textit{time behaviour} de la norma
ISO/IEC~25010:2023 \cite{iso25010}, que establece que un sistema debe responder dentro
de los tiempos que el usuario percibe como aceptables para cada tipo de operación. Para
operaciones interactivas cortas (navegación, filtros, consultas) el umbral estándar es
de 1 a 3 segundos; para operaciones de procesamiento intensivo (análisis de video) el
tiempo aceptable se determina con el usuario durante las pruebas \cite{iso25010}.

\begin{longtable}{|>{\bfseries}p{1.1cm}|p{1.8cm}|p{5.9cm}|p{4.3cm}|}
\caption{Requerimientos No Funcionales (RNF).
\label{tab:rnf}}\\
\hline
\textbf{ID} & \textbf{Categoría} & \textbf{Descripción} & \textbf{Métrica / Criterio} \\
\hline
\endfirsthead
\hline
\textbf{ID} & \textbf{Categoría} & \textbf{Descripción} & \textbf{Métrica / Criterio} \\
\hline
\endhead
\hline
\endfoot

RNF-01 & Rendimiento
       & Las operaciones de navegación e interacción con la interfaz (cargar el
         dashboard, filtrar experimentos, abrir resultados) deben responder en menos
         de 3~s. Derivado del atributo \textit{time behaviour} de la ISO/IEC~25010:2023
         \cite{iso25010}.
       & Tiempo de respuesta $\leq 3$~s en el 95\,\% de las solicitudes en condiciones
         normales de red. \\
\hline
RNF-02 & Rendimiento
       & \textbf{[PENDIENTE — entrevista con el Dr.\ Sandino]:} El tiempo máximo
         aceptable para analizar un video de 5~min (Día~2, resolución típica de celular)
         se definirá con el laboratorio durante la fase de pruebas y validación.
       & Tiempo máximo a medir y documentar durante las pruebas con videos reales
         del laboratorio. \\
\hline
RNF-03 & Disponibilidad
       & El sistema debe estar disponible en el horario de trabajo del laboratorio.
         No se requiere disponibilidad 24/7 para el prototipo.
       & 95\,\% de disponibilidad de lunes a viernes, 8:00--20:00~h. \\
\hline
RNF-04 & Seguridad
       & Solo los usuarios autenticados pueden acceder al sistema. La autenticación
         usa tokens JWT con tiempo de expiración configurable. Los datos de cada
         investigador no son visibles para otros usuarios sin autorización explícita.
       & Autenticación con JWT. Datos aislados por usuario en la base de datos.
         Verificado en las pruebas de integración de la fase de implementación. \\
\hline
RNF-05 & Seguridad
       & Las contraseñas deben almacenarse cifradas con sal, nunca en texto plano.
         Las contraseñas temporales asignadas por el Admin se marcan para cambio
         obligatorio al primer ingreso.
       & Hash con sal usando bcrypt. Política de contraseña temporal verificada
         en prueba funcional. \\
\hline
RNF-06 & Usabilidad
       & Un investigador sin experiencia en programación debe poder crear un
         experimento, subir un video, esperar el análisis y descargar el reporte
         sin necesitar ayuda técnica. Derivado del atributo \textit{usability} de
         la ISO/IEC~25010:2023 \cite{iso25010}.
       & Prueba con al menos un usuario real del laboratorio durante la fase de
         validación y robustez. Tasa de completación de tarea $\geq$ 90\% sin
         asistencia. \\
\hline
RNF-07 & Compatibilidad
       & El sistema debe funcionar desde Chrome, Firefox y Edge en sus versiones
         actuales, sin instalar nada en el equipo del usuario.
       & Pruebas en los tres navegadores en sus versiones actuales al cierre de
         la fase de implementación de la interfaz web. \\
\hline
RNF-08 & Almacenamiento
       & Los videos se borran automáticamente 30~días después de que el análisis
         termina exitosamente. Los resultados y reportes se conservan
         indefinidamente. El sistema alerta al usuario cuando un video está a menos
         de 7~días de ser borrado.
       & Tarea programada (\textit{cron job}) diaria. Aviso visible en dashboard
         con indicador de días restantes. \\
\hline
RNF-09 & Almacenamiento
       & El Administrador recibe una alerta en el panel cuando el uso del disco
         supera el 80\,\%. A los 90\,\%, el sistema elimina automáticamente los
         videos más antiguos, notificando al investigador correspondiente.
       & Umbral de alerta: 80\,\%. Umbral de limpieza automática: 90\,\%.
         Los valores 80/90\,\% son umbrales estándar de gestión de capacidad
         en sistemas de almacenamiento \cite{pressman2014}; dejan margen
         operativo antes de que el disco se llene y evitan pérdida de datos
         por falta de espacio. \\
\hline
RNF-10 & Portabilidad
       & El sistema debe poder levantarse con Docker Compose en cualquier equipo
         con Docker instalado, sin dependencias adicionales en el sistema anfitrión.
         Derivado del proceso de despliegue de la ISO/IEC/IEEE~12207:2017
         \cite{iso12207}.
       & Despliegue documentado y verificado en al menos dos entornos distintos. \\
\hline
RNF-11 & Escalabilidad
       & El módulo de clasificación debe poder reconfigurarse para soportar
         paradigmas conductuales distintos al FST sin modificar la plataforma web.
       & Diseño modular verificado en la revisión de arquitectura al cierre
         del bloque de diseño. \\
\hline
RNF-12 & Mantenibilidad
       & El código debe estar documentado y separado en módulos independientes
         (backend, pipeline de IA, frontend), siguiendo PEP~8 en Python y ESLint
         en JavaScript.
       & Revisión de calidad de código en cada cierre de bloque de
         implementación. \\
\hline
RNF-13 & Concurrencia
       & \textbf{[PENDIENTE — entrevista con el Dr.\ Sandino]:} Número de
         usuarios y análisis simultáneos que debe soportar el sistema, e
         infraestructura disponible en el laboratorio. Definir antes de la
         fase de pruebas de carga.
       & Pendiente. Se documenta antes de iniciar las pruebas de carga. \\
\hline

\end{longtable}

% -----------------------------------------------------------
\subsection{Reglas de negocio}
% -----------------------------------------------------------

Las reglas de negocio son los límites que el sistema siempre debe respetar,
independientemente de cómo esté implementado. Algunas vienen del protocolo
experimental del FST, otras del diseño de la interfaz y otras del contexto
institucional del laboratorio \cite{pressman2014}. Los requerimientos que dependen
de estas reglas las referencian explícitamente.

\begin{longtable}{|>{\bfseries}p{1.1cm}|p{8.1cm}|p{3.9cm}|}
\caption{Reglas de negocio.
\label{tab:rn}}\\
\hline
\textbf{ID} & \textbf{Regla} & \textbf{Justificación / RF relacionados} \\
\hline
\endfirsthead
\hline
\textbf{ID} & \textbf{Regla} & \textbf{Justificación / RF relacionados} \\
\hline
\endhead
\hline
\endfoot

RN-01 & Las credenciales incorrectas muestran un mensaje de error genérico que no
        especifica si el correo o la contraseña es lo que falló.
      & Seguridad básica contra enumeración de usuarios. Ver RF-01. \\
\hline
RN-02 & Un video pertenece a exactamente un experimento. No se puede reutilizar un
        video ya subido en un experimento distinto.
      & Trazabilidad de resultados. Ver RF-09. \\
\hline
RN-03 & Un experimento tiene como máximo dos videos: uno del Día~1 (20~min) y uno
        del Día~2 (5~min).
      & Así funciona el protocolo FST del laboratorio \cite{slattery2012}. Ver RF-09. \\
\hline
RN-04 & El análisis lo inicia el sistema automáticamente al terminar la carga. El
        investigador no puede iniciarlo, pausarlo, cancelarlo ni reiniciarlo desde
        la interfaz.
      & Evita estados inconsistentes en la cola. Ver RF-12, RF-16. \\
\hline
RN-05 & Los videos se borran automáticamente 30~días después de que el análisis
        termina exitosamente. La acción es irreversible. El sistema avisa con al
        menos 7~días de anticipación.
      & Control de almacenamiento. Ver RF-24, RF-23, RNF-08. \\
\hline
RN-06 & Los resultados y reportes (PDF y CSV) se conservan indefinidamente aunque
        el video original ya haya sido borrado.
      & Los datos de investigación tienen valor a largo plazo. Ver RF-25. \\
\hline
RN-07 & En cada cuadro del video, a cada animal se le asigna exactamente una
        conducta: nado activo, inmovilidad o escalamiento. No pueden coexistir dos
        conductas al mismo tiempo para el mismo animal.
      & El FST define estas conductas como mutuamente excluyentes
        \cite{armario2021,cryan2005}. Ver RF-17, RF-18. \\
\hline
RN-08 & Cada investigador solo puede ver sus propios experimentos y resultados. El
        Administrador puede ver y gestionar los de cualquier investigador.
      & Control de acceso. Ver RF-04, RF-29. \\
\hline
RN-09 & Borrar un experimento es permanente e irreversible: se eliminan el video,
        los resultados y los reportes. El sistema requiere confirmación explícita
        antes de ejecutarlo.
      & Previene borrados accidentales. Ver RF-26. \\
\hline
RN-10 & El sistema rechaza cualquier archivo que no sea \texttt{.mp4} antes de
        intentar subirlo al servidor. La validación de formato ocurre en el cliente.
      & Evita consumir recursos con archivos incompatibles. Ver RF-08, RF-10. \\
\hline
RN-11 & La detección de cilindros requiere una confianza mínima de 0.70. Si no se
        alcanza ese umbral, el pipeline se detiene y genera un reporte de
        diagnóstico descargable automáticamente. El umbral de 0.70 es un valor
        conservador estándar en detección con YOLOv8 para fondos consistentes
        \cite{stadler2024_yolov8}; se ajustará con base en las pruebas con los
        videos reales del laboratorio durante la fase de validación.
      & Evita que el sistema produzca resultados incorrectos sin advertir al
        usuario. Ver RF-14, RF-16. \\
\hline
RN-12 & La comparación Día~1 vs.\ Día~2 solo se genera si ambos videos del mismo
        experimento fueron procesados exitosamente.
      & No tiene sentido comparar con datos incompletos. Ver RF-20. \\
\hline

\end{longtable}

% -----------------------------------------------------------
\section{Análisis de las herramientas a utilizar}
% -----------------------------------------------------------

La elección del \textit{stack} tecnológico se hizo evaluando madurez, compatibilidad,
disponibilidad de documentación y experiencia previa del equipo. Las herramientas se
agrupan en cuatro áreas.

\subsection{Detección de objetos y seguimiento}

\begin{longtable}{|p{2.3cm}|p{5.4cm}|p{5.4cm}|}
\caption{Técnicas de detección de objetos y seguimiento.
\label{tab:herr_cv}}\\
\hline
\textbf{Herramienta} & \textbf{Características principales} & \textbf{Por qué se eligió} \\
\hline
\endfirsthead
\hline
\textbf{Herramienta} & \textbf{Características principales} & \textbf{Por qué se eligió} \\
\hline
\endhead
\hline
\endfoot

OpenCV 4.x \cite{bradski2000}
    & Biblioteca estándar para procesamiento de imágenes y video en Python. Incluye
      sustracción de fondo, umbralización de Otsu, morfología, contornos y los
      trackers CSRT y KCF.
    & Es la herramienta de referencia para este tipo de trabajo y es compatible con
      YOLOv8 y TensorFlow sin configuración extra. \\
\hline
Ultralytics YOLOv8 \cite{jocher2023}
    & Detector de objetos estado del arte. Con 100--350 imágenes etiquetadas logra
      detección casi perfecta en fondos consistentes \cite{stadler2024_yolov8}.
    & El detector de rata ya fue entrenado (\texttt{weights/rat.pt}) y su viabilidad
      técnica está confirmada con los videos del laboratorio. \\
\hline
Sustracción de fondo por mediana \cite{piccardi2004}
    & El fondo se estima como la mediana píxel a píxel de los primeros $N$ cuadros,
      evitando que una rata inmóvil al inicio contamine el modelo.
    & Más simple y predecible que el GMM \cite{zivkovic2004} para el caso de cámara
      fija. No requiere parámetros de color ni modelos adaptativos, lo que reduce
      la sensibilidad a variaciones de iluminación entre sesiones. \\
\hline
Tracker CSRT \cite{lukezic2017}
    & Tracker de correlación con mapa de confiabilidad espacial. Más preciso que
      KCF aunque más lento.
    & Primera opción de respaldo cuando YOLO no detecta a la rata en un cuadro.
      KCF \cite{henriques2015} como respaldo secundario. \\
\hline

\end{longtable}

\subsection{Aprendizaje profundo para clasificación de conductas}

\begin{longtable}{|p{2.3cm}|p{5.4cm}|p{5.4cm}|}
\caption{Arquitecturas de aprendizaje profundo para clasificación de conductas.
\label{tab:herr_dl}}\\
\hline
\textbf{Herramienta / Arquitectura} & \textbf{Características principales} & \textbf{Por qué se eligió} \\
\hline
\endfirsthead
\hline
\textbf{Herramienta / Arquitectura} & \textbf{Características principales} & \textbf{Por qué se eligió} \\
\hline
\endhead
\hline
\endfoot

ResNet-18 \cite{he2016}
    & Red con 18~capas y conexiones residuales, preentrenada en ImageNet. Ligera
      ($< 50$~MB) y rápida de entrenar con datasets pequeños via transfer learning
      \cite{pan2010}.
    & Primera opción para el clasificador de conducta. Su viabilidad se
      confirmará durante la fase de implementación del clasificador. \\
\hline
ResNet-50 \cite{he2016}
    & Versión más profunda (50~capas). Mayor capacidad que ResNet-18 a costa de más
      tiempo de entrenamiento e inferencia.
    & Se usará si ResNet-18 no alcanza $\kappa > 0.80$. \\
\hline
TensorFlow 2.x / Keras \cite{abadi2016}
    & Framework de aprendizaje profundo con API de alto nivel para fine-tuning de
      modelos preentrenados.
    & Ya se usa en el proyecto y es compatible con Docker para el despliegue. \\
\hline
scikit-learn \cite{pedregosa2011}
    & Biblioteca de ML para métricas y validación cruzada.
    & Se usa para calcular el coeficiente $\kappa$ de Cohen, F1 por clase y MAE. \\
\hline
BORIS \cite{friard2016}
    & Software de anotación conductual cuadro a cuadro. Exporta en formatos
      compatibles con Python.
    & Los autores del proyecto anotan los cuadros de video con BORIS usando los
      registros conductuales del laboratorio como referencia, generando el dataset
      de entrenamiento y el gold standard de validación. \\
\hline

\end{longtable}

\subsection{Desarrollo web e infraestructura}

\begin{longtable}{|p{2.3cm}|p{5.4cm}|p{5.4cm}|}
\caption{Herramientas de desarrollo web e infraestructura.
\label{tab:herr_web}}\\
\hline
\textbf{Herramienta} & \textbf{Características principales} & \textbf{Por qué se eligió} \\
\hline
\endfirsthead
\hline
\textbf{Herramienta} & \textbf{Características principales} & \textbf{Por qué se eligió} \\
\hline
\endhead
\hline
\endfoot

Flask (Python) \cite{grinberg2018}
    & Microframework ligero para APIs REST. Extensiones para JWT
      (Flask-JWT-Extended), ORM (SQLAlchemy) y tareas asíncronas.
    & Suficiente para el tamaño del prototipo. Integra el pipeline de IA sin
      necesitar un framework más pesado. \\
\hline
React.js + Vite
    & Framework frontend basado en componentes. Vite reduce los tiempos de
      compilación durante el desarrollo.
    & Adecuado para la barra de progreso en tiempo real y la interfaz reactiva
      que requiere el sistema. \\
\hline
PostgreSQL
    & Base de datos relacional con transacciones ACID y código abierto.
    & Cubre los patrones de acceso del sistema y se integra con SQLAlchemy. \\
\hline
Docker + Docker Compose \cite{merkel2014}
    & Empaqueta todos los servicios (backend, frontend, base de datos, worker) en
      contenedores y los levanta con un solo comando.
    & Cumple el atributo de portabilidad de la ISO/IEC/IEEE~12207:2017
      \cite{iso12207}. El laboratorio no instala nada más allá de Docker. \\
\hline

\end{longtable}

% -----------------------------------------------------------
\section{Análisis de riesgos}
% -----------------------------------------------------------

La identificación de riesgos sigue el modelo de la norma ISO~31000 \cite{iso31000}, que
define el riesgo como la combinación entre la probabilidad de un evento y su impacto.
Para cada riesgo se documenta una acción concreta de mitigación.

\begin{longtable}{|>{\bfseries}p{0.9cm}|p{3.4cm}|p{1.3cm}|p{1.2cm}|p{1.1cm}|p{4.6cm}|}
\caption{Matriz de riesgos del proyecto.
\label{tab:riesgos}}\\
\hline
\textbf{ID} & \textbf{Descripción} & \textbf{Prob.} & \textbf{Impacto} & \textbf{Nivel} & \textbf{Mitigación} \\
\hline
\endfirsthead
\hline
\textbf{ID} & \textbf{Descripción} & \textbf{Prob.} & \textbf{Impacto} & \textbf{Nivel} & \textbf{Mitigación} \\
\hline
\endhead
\hline
\endfoot

R-01 & Videos con reflejos, iluminación variable o ratas de pelaje claro que hacen
      fallar el pipeline (confianza de detección $< 0.70$, ver RN-11).
    & Alta & Alto & Alto
    & Definir requisitos mínimos de calidad de video con el laboratorio. El sistema
      aplica corrección CLAHE (\textit{Contrast Limited Adaptive Histogram
      Equalization}) \cite{bradski2000} como primer intento automático de mejora
      de contraste antes de abortar el pipeline; si después de CLAHE la confianza
      sigue por debajo de 0.70, se genera el reporte de diagnóstico (ver RN-11). \\
\hline
R-02 & El dataset de cuadros etiquetados no alcanza para entrenar un clasificador
      con $\kappa > 0.80$.
    & Alta & Alto & Alto
    & Ampliar el dataset anotando más cuadros con BORIS. Usar data augmentation.
      Si escalamiento no se distingue con precisión, reportarlo junto con nado
      como ``conducta activa'' (ver alcance, cap.~1). \\
\hline
R-03 & El laboratorio no proporciona los registros conductuales en el tiempo
      necesario para el entrenamiento y validación del clasificador.
    & Media & Alto & Medio
    & Acordar un calendario de entrega de registros con el laboratorio al inicio
      de la fase de implementación y preparar los materiales de anotación con
      anticipación para aprovechar cada sesión con el Dr.\ Sandino. \\
\hline
R-04 & El tiempo de análisis de un video de 5~min supera lo que el laboratorio
      puede asumir en su flujo de trabajo.
    & Media & Medio & Medio
    & Medir tiempos por etapa del pipeline. Evaluar reducción de resolución o
      uso de GPU. El umbral aceptable se define con el laboratorio durante la
      fase de pruebas (ver RNF-02). \\
\hline
R-05 & La integración del módulo de IA con el backend via cola asíncrona resulta
      más compleja de lo esperado.
    & Media & Medio & Medio
    & Definir contratos de API al inicio de la fase de implementación y hacer
      una prueba de integración de extremo a extremo antes de agregar más
      funcionalidad. \\
\hline
R-06 & El disco del servidor se llena por acumulación de videos durante los 30~días
      de retención.
    & Baja & Medio & Bajo
    & El sistema incluye alertas al 80\,\% de uso de disco y limpieza automática
      al 90\,\% (ver RF-28, RNF-09). \\
\hline
R-07 & El laboratorio solicita cambios en los requerimientos a mitad del desarrollo.
    & Baja & Alto & Medio
    & Diseño modular del clasificador para agregar clases. Gestionar cambios
      dentro del proceso de revisión de sprint. \\
\hline
R-08 & El escalamiento y el nado activo se ven tan parecidos desde la vista cenital
      que el modelo no los distingue con suficiente precisión.
    & Alta & Medio & Medio
    & Contemplado en el alcance del cap.~1: ambas conductas se reportan juntas
      como ``conducta activa'' si no se alcanza la precisión objetivo. \\
\hline

\end{longtable}

% -----------------------------------------------------------
\section{Análisis de factibilidad}
% -----------------------------------------------------------

\subsection{Factibilidad tecnológica}

La factibilidad tecnológica consiste en evaluar si el proyecto puede desarrollarse con
la infraestructura técnica, el conocimiento especializado y los recursos disponibles.
El objetivo es determinar si las tecnologías actuales y el personal son suficientes para
implementar el prototipo de forma eficiente y sostenible \cite{pressman2014}.

\subsubsection{Evaluación de recursos actuales y brechas tecnológicas}

El equipo de desarrollo cuenta con experiencia en las herramientas del stack tecnológico
seleccionado ---Python, OpenCV, TensorFlow, Flask y React--- adquirida en cursos de la
carrera en la ESCOM. Todas las herramientas son de código abierto, están bien
documentadas y tienen comunidades activas que reducen el riesgo de incompatibilidades.

La brecha tecnológica más relevante identificada es el dataset de cuadros etiquetados
para entrenar el clasificador de conducta. Con el número de videos disponibles actualmente
en el laboratorio, ese dataset requiere un esfuerzo de anotación manual con BORIS que
es viable pero no trivial. La estrategia de transfer learning con ResNet-18 reduce ese
requerimiento a unos pocos cientos de cuadros por clase, lo que está dentro del alcance
del proyecto \cite{pan2010,meng2024_dsaan}. Para el entrenamiento del modelo, si el
laboratorio no dispone de GPU, Google Colab ofrece hasta 12~horas de GPU por sesión de
forma gratuita, lo que es suficiente para el fine-tuning de ResNet-18.

En cuanto al hardware para el despliegue, el sistema no requiere infraestructura
especializada: cualquier equipo con Docker instalado puede correr el prototipo, lo que
elimina dependencias de hardware específico en el laboratorio.

\subsubsection{Simulación, modelado y pruebas}

Antes de la implementación completa del sistema, se realizaron pruebas preliminares de
detección de objetos con YOLOv8 usando un dataset inicial de los videos del laboratorio.
Los resultados de esa prueba confirman que el enfoque es viable con las condiciones de
grabación reales del laboratorio: vista cenital, cuatro cilindros simultáneos, cámara de
celular \cite{stadler2024_yolov8}. Esa prueba reduce el riesgo técnico de la etapa de
detección antes de comprometer recursos en el desarrollo completo.

Para el clasificador de conducta, la literatura reporta precisiones superiores al
88\,\% con datasets de menos de 100 imágenes cuando se usa transfer learning desde
ImageNet \cite{meng2024_dsaan}, lo que da respaldo empírico a la viabilidad de esa
etapa con el dataset disponible.

\subsubsection{Evaluación de compatibilidad e integración}

El prototipo tiene una arquitectura modular que separa el pipeline de IA (Python,
OpenCV, TensorFlow) de la plataforma web (Flask, React, PostgreSQL) mediante una
cola de tareas asíncrona. Esa separación permite desarrollar y probar cada módulo de
forma independiente y facilita la integración progresiva conforme avanza el desarrollo.

La contenedorización con Docker garantiza que el sistema se comporta igual en el equipo
de desarrollo, en el laboratorio y en cualquier servidor institucional, lo que elimina
la clase más común de problemas de compatibilidad en el despliegue \cite{merkel2014}.

\subsection{Factibilidad operativa}

El objetivo de la factibilidad operativa es establecer si el sistema propuesto puede
emplearse adecuadamente en el contexto académico y experimental para el que fue
concebido. Este proyecto responde a una necesidad concreta expresada por el
Dr.\ César Augusto Sandino Reyes López, investigador del Laboratorio de Neurociencia
Conductual de la Escuela Nacional de Medicina y Homeopatía (ENMyH) del IPN, quien
realiza experimentos de nado forzado de forma regular y actualmente analiza los videos
de forma manual.

Para respaldar esta factibilidad, se realizan reuniones con el Dr.\ Sandino en las
que se aplica un instrumento de recolección de datos ---una entrevista estructurada
con preguntas abiertas--- orientado a obtener información concreta sobre la demanda
del sistema, los indicadores más relevantes para el laboratorio, las condiciones reales
de operación y los criterios de aceptación desde la perspectiva del usuario experto.
Los resultados de esa entrevista se incorporan a este documento antes de la entrega
final; el instrumento completo se incluye en los anexos.

\textbf{[PENDIENTE --- entrevista con el Dr.\ Sandino]:}

Con base en las reuniones de arranque del proyecto y en el análisis del protocolo
experimental del laboratorio (cap.~1), se pueden anticipar las siguientes ventajas
operativas del sistema: reducción del tiempo de análisis de varias horas a minutos,
estandarización del criterio de evaluación entre sesiones y evaluadores, generación
automática de reportes en los formatos que ya usa el laboratorio (PDF, CSV) y
posibilidad de comparar grupos y tratamientos desde una interfaz centralizada.

\subsection{Factibilidad económica}

El proyecto no requiere licencias de software. Todas las herramientas del stack son
de código abierto con licencias compatibles con uso académico. El desarrollo lo realiza
el equipo del TT sin costo adicional, y la infraestructura de despliegue usa el equipo
del laboratorio o las computadoras del equipo de desarrollo.

\begin{longtable}{|p{5.8cm}|p{7.5cm}|}
\caption{Análisis de costos del proyecto.
\label{tab:costos}}\\
\hline
\textbf{Concepto} & \textbf{Costo estimado} \\
\hline
\endfirsthead
\hline
\textbf{Concepto} & \textbf{Costo estimado} \\
\hline
\endhead
\hline
\endfoot

Licencias de software (Python, Flask, React, PostgreSQL, Docker, YOLOv8, OpenCV, BORIS)
    & \$0.00 --- todas de código abierto con licencias compatibles con uso académico. \\
\hline
Entrenamiento del modelo (GPU)
    & \$0.00 --- Google Colab gratuito (hasta 12~h de GPU por sesión). \\
\hline
Infraestructura de despliegue
    & Sin costo adicional --- equipo del laboratorio o computadora del equipo de TT. \\
\hline
Almacenamiento de videos
    & $\approx$500~MB por experimento de 5~min en MP4. La política de eliminación
      automática a 30~días (RN-05) evita que el espacio crezca sin control. \\
\hline
\textbf{Costo total de herramientas y desarrollo} & \textbf{\$0.00} \\
\hline

\end{longtable}

El proyecto no requiere financiamiento externo ni aprobaciones institucionales para
iniciar el desarrollo.

\subsection{Factibilidad legal}

Todas las herramientas del proyecto usan licencias de código abierto que permiten el
uso académico sin restricciones: PSF (Python), BSD (Flask), MIT (React), PostgreSQL
license (PostgreSQL), Apache~2.0 (Docker, OpenCV), AGPL-3.0 para uso no comercial
(YOLOv8) y GNU GPL v3 (BORIS \cite{friard2016}).

Los videos son propiedad del grupo de investigación de la ENMyH-IPN y su uso está
respaldado por la colaboración formal entre los directores del TT y el laboratorio.
El equipo de desarrollo no interviene en ningún experimento con animales; eso lo hace
el laboratorio bajo la NOM-062-ZOO-1999 \cite{nom062}. El desarrollo del sistema
cumple con los Lineamientos de Propiedad Intelectual y Transferencia de Conocimiento
del IPN (2023).

% -----------------------------------------------------------
\section{Análisis de sostenibilidad}
% -----------------------------------------------------------

\subsection{Sostenibilidad ambiental}

Reducir el análisis manual quita la necesidad de que los investigadores vayan al
laboratorio solo para cronometrar videos. Al ser web, el sistema se puede usar desde
cualquier lugar. La política de borrar videos a los 30~días evita que el disco crezca
sin control. Según las directrices de la norma ISO~14064-1 \cite{iso14064}, el consumo
energético de procesar un video de 5~min en una computadora local es marginal comparado
con las horas de trabajo que reemplaza.

\subsection{Sostenibilidad social e institucional}

El sistema está hecho para el laboratorio de la ENMyH, pero su diseño modular permite
adaptarlo a otros paradigmas conductuales (laberinto en cruz elevada, campo abierto,
etc.) reconfigurando solo el módulo de clasificación. Eso abre la puerta a que otros
grupos del IPN lo usen sin desarrollarlo desde cero.

Si funciona bien, el Dr.\ Sandino deja de cronometrar videos a mano. Eso es lo que
el laboratorio necesita; cualquier beneficio más amplio ---reproducibilidad entre
grupos, escalabilidad a otros paradigmas--- depende de que primero funcione para
el caso para el que fue construido.

\subsection{Mantenibilidad del sistema}

El sistema se entrega como prototipo funcional. Para mantenerlo hace falta conocimiento
básico de Python y Docker. La documentación técnica de este documento y el
\texttt{README} del repositorio están diseñados para que cualquiera que retome el
proyecto pueda entenderlo sin el equipo original. Las herramientas principales (YOLOv8,
Flask, React, PostgreSQL) tienen desarrollo activo y comunidades grandes.